\documentclass[12pt]{article}
\usepackage[utf8]{inputenc}
\usepackage[T1]{fontenc}
\usepackage[margin=1in]{geometry}
\usepackage{setspace}
\usepackage{xcolor}
\usepackage[colorlinks=true, urlcolor=blue!60!black, linkcolor=blue!60!black]{hyperref}
\setlength{\parindent}{0em}
\setlength{\parskip}{0.45em}
\renewcommand{\baselinestretch}{1.0}
\begin{document}
\begin{center}
{\Large \textit{McCulloch v. Maryland}, 17 U.S. (4 Wheat.) 316 (1819)\par}
\vspace{0.3em}
{\normalsize Case Brief --- One-Page Edition\par}
\end{center}

\textbf{Memory Jogger.} Implied powers and intergovernmental tax immunity: the national bank case (\href{https://supreme.justia.com/cases/federal/us/17/316/}{17 U.S. 316}).

\textbf{Facts.} Congress chartered the Second Bank of the United States in 1816. Maryland taxed banks not chartered by its legislature, in practice only the Bank's Baltimore branch. Cashier James McCulloch issued notes without paying and was sued for the penalties.

\textbf{Procedural History.} Maryland's courts ruled against McCulloch; he sought review by writ of error.

\textbf{Votes.} Reversed, unanimously. Marshall wrote for the Court; there were no separate opinions.

\textbf{Holding.} Congress may use any appropriate, unprohibited means plainly adapted to its enumerated powers, so it may charter a bank. States may not tax or burden federal operations.

\textbf{Reasoning.} The Constitution comes from the people and marks only great outlines; the Tenth Amendment omits ``expressly.'' ``Necessary'' means useful or conducive, as ``absolutely necessary'' in Article I, Section 10 shows, and a strict reading would cripple the government. Because the power to tax involves the power to destroy, one state may not tax an instrument of the whole nation, whose people it does not represent.

\textbf{Examples.} \emph{Same-Side:} federal civil commitment of dangerous prisoners. \emph{Opposite-Side:} a general state property tax on a federal bank's building. \emph{Fence-Sitter:} a remote chain of implied powers.

\textbf{Critique.} Originalists say Marshall loosened ``necessary and proper'' too far; nationalists say he conceded too much to enumeration.

\textbf{Key Quotation.} ``Let the end be legitimate, let it be within the scope of the Constitution, and all means which are appropriate, which are plainly adapted to that end, which are not prohibited, but consist with the letter and spirit of the Constitution, are Constitutional.'' --- Chief Justice Marshall, Opinion of the Court, 17 U.S. (4 Wheat.) at 421.

\end{document}
